% Lean module: MonadFormal.Temporal.Variational.Regularity
\subsubsection{Regularity from the weak equation}
\label{sec:variational-regularity}

\begin{lemma}[Mean-zero orthogonality and integration by parts]
\label[lemma]{lem:variational-orthogonality}
If $k\in V$ is orthogonal to every $h\in H$, then $k$ is
constant almost everywhere. For $g\in V$ and $h\in H$,
\[
 \langle Pg,h\rangle_{L^2(\mu)}
 =-\int g(t)(Ph)(t)\,d\mu(t).
\]
Here the bounded continuous function $Pg$ is also regarded as an
element of $V$.
\end{lemma}
\begin{proof}
Put $m=\int k\,d\mu$. Since $\mu([0,1])=1$, the function
$h=k-m$ belongs to $H$. Constants are orthogonal to $H$, so
$\lVert h\rVert_2^2=\langle k-m,h\rangle=0$ and $k=m$
almost everywhere. For the second identity, $Pg$ and $Ph$ are
absolutely continuous, with derivatives $g$ and $h$ almost
everywhere. Integration by parts has zero boundary term because
$Ph(0)=Ph(1)=0$.
\leandef{MonadFormal.Temporal.Variational.Interval.primitiveLp}
\leanproof{MonadFormal.Temporal.Variational.Interval.endpoint_const,MonadFormal.Temporal.Variational.Interval.const_inner_variation,MonadFormal.Temporal.Variational.Interval.eq_const_of_orthogonal,MonadFormal.Temporal.Variational.Interval.primitiveLp_ae,MonadFormal.Temporal.Variational.Interval.primitive_pairing}
\end{proof}

\begin{theorem}[Weak equation determines an absolutely continuous velocity]
\label{thm:variational-regularity}
Let $f\in H$, $g\in V$ and $\lambda\in\mathbb R$ satisfy
\[
 \langle f,h\rangle+\lambda\int g(t)(Ph)(t)\,d\mu(t)=0
 \qquad(h\in H).
\]
There is a constant $c$ such that $f=c+\lambda Pg$ almost
everywhere. Consequently, $x_f$ has a continuous velocity
\[
 v(t)=b-a+c+\lambda(Pg)(t),
\]
which is absolutely continuous on $[0,1]$. The equality
$x_f'(t)=v(t)$ holds at every interior point, and
$v'=x_f''=\lambda g$ almost everywhere on $[0,1]$.
\end{theorem}
\begin{proof}
By \Cref{lem:variational-orthogonality}, the weak equation says
that $f-\lambda Pg$ is orthogonal to $H$, hence is constant.
The primitive $Pg$ is continuous and absolutely continuous, so
the displayed formula gives the same regularity for $v$.
Integrating the continuous representative $c+\lambda Pg$ of $f$
upgrades the derivative of $Pf$ to every interior point by the
fundamental theorem of calculus. Finally $(Pg)'=g$ almost
everywhere. Since $x_f'=v$ on the open interval, their derivatives
also agree wherever $v'$ exists there.
\leanproof{MonadFormal.Temporal.Variational.Interval.weak_velocity_representation,MonadFormal.Temporal.Variational.Interval.primitive_hasDerivAt_of_continuous,MonadFormal.Temporal.Variational.Interval.weak_path_regularity,MonadFormal.Temporal.Variational.Interval.second_derivative_ae}
\end{proof}
